% !TeX root = ../../Book3.tex
% 纲要标题：### B.3 分解、参数与组合

\section{分解、参数与组合}
\label{b3:sec:B03}

一个分解或正规化的候选输出，必须同时满足包含关系、终止条件与所声称的结构性质。参数化方程还要求说明输出在哪些参数取值下成立，不能把一次泛型计算当作全部纤维的答案。

% 纲要细目（与计划纲要同步）：
% * 准素分解、根理想、饱和与正规化的有效问题
% * 参数化方程的综合 Gröbner 系统、正则链及参数空间分层
% * 一般参数值的答案不代表所有特殊化；退化参数需要独立分支
% * 单项式理想、组合结构与 Hilbert 数据；分解结果的唯一性条件
%
% #### B.3.1 准素分解、根理想、饱和与正规化的有效问题
% #### B.3.2 综合 Gröbner 系统、正则链与参数空间分层
% #### B.3.3 一般参数、特殊化与退化分支
% #### B.3.4 单项式理想、组合结构与分解的非唯一性


\subsection{准素分解、根理想、饱和与正规化的有效问题}
\label{b3:subsec:B03:01}

正文已经证明，饱和、根理想、准素分解和整闭包具有各自的代数意义。有效计算的任务是从有限输入得到有限输出，并附上足以识别该输出的依据。四种输出不能混为一谈：根理想舍去幂零结构；准素分解保存理想本身；饱和选择某些分支；正规化则通常改变坐标环而非只把原理想换成另一个根理想。

\begin{example}\label{b3:exm:B-decomposition}
在 \(k[x,y]\) 中，
\[
I=(x^2,xy)=(x)\cap(x^2,y),\qquad
\sqrt I=(x),\qquad I:y^\infty=(x).
\]
交式可逐项检查：若 \(xh\in(x^2,y)\)，模 \(y\) 后 \(h(x,0)\) 被 \(x\) 整除，所以 \(h\in(x,y)\)。第二个理想的商为 \(k[x]/(x^2)\)，故它是 \((x,y)\)-准素的。饱和的等式来自 \(yx\in I\) 及 \(I\subset(x)\)，其中 \(y\notin(x)\)。这个例子中的饱和移除了嵌入的加厚结构，而根理想同时忘记了全部幂零信息。
\end{example}

对于域上的有限多项式输入，交和饱和可用\cref{b3:prop:ideal-intersection-elimination,b3:prop:colon-saturation-elimination}的辅助变量构造化为消去。核验一个候选准素分解时，还要证明各分量准素、交等于原理想，以及在声称不可约简时没有冗余。仅列出若干零点集的并不够，因为它不能区分 \(I\) 与 \(\sqrt I\)。

\ProseParagraphBreak
正规化的输出则应包含有限代数及其到总分式环的嵌入。例如 \(A=k[t^2,t^3]\subset k[t]\) 的正规化为 \(k[t]\)：\(t\) 满足首一方程 \(T^2-t^2=0\)，且 \(k[t]=A+At\)；两者分式域相同，而 \(k[t]\) 是整闭环。若一个分式对 \(A\) 整，则同一首一方程也说明它对 \(k[t]\) 整，因此它属于 \(k[t]\)。这完整核验了该例，但一般正规化算法还须处理有限性假设、分支及系数域，本附录不以此例代替它们。

\subsection{综合 Gröbner 系统、正则链与参数空间分层}
\label{b3:subsec:B03:02}

参数计算的难处不是增加几个字母，而是某些系数随参数取值变为零。把参数当作有理函数域的元素，可得到一般情形的答案；要覆盖全部参数，还须记录除法中要求哪些系数非零。

\begin{definition}\label{b3:def:B-comprehensive-system}
设 \(k\) 为域，\(\overline k\) 为其代数闭包，\(a=(a_1,\ldots,a_r)\) 为参数，\(x=(x_1,\ldots,x_n)\) 为未知量，\(F\subset k[a,x]\) 为有限集，并固定 \(x\) 的全局项序。若有限个由参数多项式等式和非等式定义的集合 \(C_\nu\subset\overline{k}^{\,r}\) 覆盖参数空间，并各附一列可在 \(C_\nu\) 上有定义的多项式 \(G_\nu\)，使对每个 \(c\in C_\nu\)，\(G_\nu(c,x)\) 是 \((F(c,x))\) 的 Gröbner 基，则这些分支数据称为一个\term{综合 Gröbner 系统}[comprehensive Gröbner system]。
\end{definition}

这一定义描述输出应满足什么，不声称任意一次一般参数计算已经给出这样的覆盖。若分支重叠，其特殊化答案须描述同一输入理想；若要求互不相交，还需进一步细分条件。

\ProseParagraphBreak
三角形方程组提供另一种组织方式。按变量 \(x_1<\cdots<x_n\) 排列，一个非常数多项式实际出现的最高变量称为它的主变量；把它视为主变量的多项式时，其首系数称为初式。

\begin{definition}\label{b3:def:B-regular-chain}
设 \(K\) 为域，\(T=(t_1,\ldots,t_s)\subset K[x_1,\ldots,x_n]\) 为非常数多项式列，主变量严格递增，初式依次为 \(h_1,\ldots,h_s\)。令 \(H_{i-1}=h_1\cdots h_{i-1}\)，并约定空乘积为 \(1\)。若每个 \(h_i\) 在
\[
K[x_1,\ldots,x_n]/
\bigl((t_1,\ldots,t_{i-1}):H_{i-1}^{\infty}\bigr)
\]
中均为非零因子，且这些商环非零，则称 \(T\) 为一条\term{正则链}[regular chain]。
\end{definition}

这个条件容许前面的商环有零因子，却排除所需初式成为零因子。例如 \((x^2,xy+1)\) 关于 \(x<y\) 是三角形方程组，但第二个初式 \(x\) 在 \(K[x,y]/(x^2)\) 中为零因子，故不是正则链。链所描述的基本集合通常是
\[
V(T)\setminus V(h),\qquad
h=\text{各初式的乘积},
\]
其闭包与饱和理想 \((T):h^\infty\) 有关。因而“正则链”既不等于任意三角生成组，也不自动等于素理想的一组生成元。一般正则链分解的构造和终止证明属于参数计算专门理论，这里用下面的分支计算说明为何需要它们。

\subsection{一般参数、特殊化与退化分支}
\label{b3:subsec:B03:03}

\begin{example}\label{b3:exm:B-parameter-branches}
取参数 \(a,b\)，未知量 \(x,y\)，方程组
\[
ax-b=0,\qquad ay=0.
\]
在参数域 \(k(a,b)\) 上得到 \(x=b/a,\ y=0\)。但完整参数空间必须分为
\[
\begin{array}{c|c|c}
\text{参数条件}&\text{特殊化理想}&\text{一个 Gröbner 基}\\ \hline
a\ne0&(x-b/a,y)&\{x-b/a,y\}\\
a=0,\ b\ne0&(1)&\{1\}\\
a=0,\ b=0&(0)&\varnothing
\end{array}
\]
三种情况分别有一个点、无解以及整个仿射平面。各行均由代入和可逆系数的除法直接核验，并且这些条件互不相交、覆盖全部参数。
\end{example}

一般参数答案没有“错误”；它准确描述开集 \(a\ne0\)。错误在于省去这个条件，把出现的分母形式上约掉以后当作全体参数的结论。\cref{b3:prop:groebner-specialization-open}给出更一般的安全情形：有限计算所用的系数条件同时不退化时，证书可以特殊化。

\ProseParagraphBreak
有些退化不改变解集的点数，却改变重数。例如 \(x^2-a\) 在代数闭域上、特征不为 \(2\) 时，\(a\ne0\) 有两个不同根，\(a=0\) 只有一个加厚的根。两个商代数的向量空间维数都为 \(2\)，所以“维数保持”不能代替根理想或可约性判定。把参数分支与商代数的乘法矩阵一起记录，可以看出简单特征值怎样合并为带幂零部分的局部代数。

\subsection{单项式理想、组合结构与分解的非唯一性}
\label{b3:subsec:B03:04}

单项式理想把成员关系化为指数向量的比较。设
\[
J=(x^{\alpha^{(1)}},\ldots,x^{\alpha^{(s)}})\subset
S=k[x_1,\ldots,x_n].
\]
单项式 \(x^\beta\) 属于 \(J\)，当且仅当某个 \(\alpha^{(i)}\) 在每个坐标上都不超过 \(\beta\)。标准单项式的计数因而可以用有限个指数区域的并来组织。

\begin{remark}\label{b3:prop:B-monomial-hilbert}
单项式理想的 Hilbert 级数计算直接使用\cref{b3:prop:hilbert-initial-count}已经证明的容斥公式：对生成元的每个子集计算最小公倍式的次数，再按子集大小交错相加，所得分子除以 \((1-t)^n\)。实际所需数据只有这些最小公倍式及其次数，不必先计算自由消解。多个子集若有相同的最小公倍式，可以合并其带符号的贡献；冗余生成元虽不影响答案，却会增加需要枚举的子集数。
\end{remark}

例如取 \(J=(x^3,x^2y,y^3)\subset k[x,y]\)。三个生成元的次数都为 \(3\)，两两最小公倍式的次数依次为 \(4,6,5\)，三个生成元的最小公倍式次数为 \(6\)。所以两个 \(t^6\) 项抵消，得到
\[
H_{k[x,y]/J}(t)
=\frac{1-3t^3+t^4+t^5}{(1-t)^2}
=1+2t+3t^2+t^3.
\]
也可以按 \(x\) 的指数分组：指数零和一时，各允许 \(y^0,y,y^2\)；指数二时只允许 \(y^0\)。由此得到七个标准单项式
\[
1,\ y,\ y^2,\ x,\ xy,\ xy^2,\ x^2,
\]
逐次计数与上面的级数一致，并给出商的长度为 \(7\)。这个例子同时核验了最小公倍式的容斥计算与直接的指数计数。

\ProseParagraphBreak
准素分解还有另一种非唯一性。前例 \((x^2,xy)\) 对每个整数 \(r\ge1\) 都有
\[
(x^2,xy)=(x)\cap(x^2,xy,y^r).
\]
第二个商环中 \(x,y\) 都幂零，且剩余域为 \(k\)，故第二分量为 \((x,y)\)-准素；交式由单项式整除即可核验。这些嵌入准素分量随 \(r\) 改变。关联素理想集合具有正文已经说明的唯一性，而具体嵌入分量并无同样唯一性。比较两个程序输出时，应比较它们是否满足所要求的分解性质，不能只因列表不同就判定其中一个错误。
